\documentclass[12pt]{article}
\usepackage{booktabs}
\usepackage{caption}
\usepackage{threeparttable}
\usepackage{dcolumn}
\usepackage{adjustbox}
\usepackage[utf8]{inputenc}
\usepackage[margin=0.5in, paperwidth=14in, paperheight=10in]{geometry}
\pagestyle{empty}
\begin{document}
\setcounter{table}{9}
\begin{table}[htbp]
\centering
\begin{threeparttable}

\begin{tabular}{lrlrr}
\toprule
Departamento & IVIG & Grupo & Productividad & PIB\\
\midrule
ATLANTICO & 0.348 & Vulnerable & 173.625 & 70.862\\
RISARALDA & 0.314 & Vulnerable & 145.678 & 26.404\\
BOLIVAR & 0.310 & Vulnerable & 162.806 & 56.924\\
CUNDINAMARCA & 0.300 & Vulnerable & 153.282 & 99.945\\
CALDAS & 0.292 & Vulnerable & 190.927 & 26.213\\
\addlinespace
TOLIMA & 0.280 & Vulnerable & 141.333 & 33.735\\
VALLE DEL CAUCA & 0.277 & Vulnerable & 152.836 & 153.564\\
CAUCA & 0.277 & Vulnerable & 150.768 & 28.537\\
BOGOTA, D.C. & 0.241 & Vulnerable & 146.220 & 395.438\\
BOYACA & 0.228 & Vulnerable & 179.730 & 42.277\\
\addlinespace
ANTIOQUIA & 0.212 & Control & 154.082 & 231.122\\
SANTANDER & 0.207 & Control & 177.290 & 101.446\\
CORDOBA & 0.195 & Control & 144.475 & 28.303\\
\bottomrule
\end{tabular}
\begin{tablenotes}
\footnotesize
\item \textit{Notas:}El PIB departamental en Billones COP
\end{tablenotes}
\end{threeparttable}
\end{table}
\end{document}
